\section{Scope, stipulated input, and attribution}\label{sec:intro}
Let $\Ecal$ be the unoriented nearest-neighbor edges of $\Z^2$. For positive
weights $w$, write $T_w(x,y)$ for the infimum of the costs of finite walks
from $x$ to $y$, counting every traversal. A ray is an infinite simple path
whose finite segments minimize cost; a bigeodesic is such a path indexed by
$\Z$. All weights in this paper are finite. A nonatomic nonnegative law
assigns probability one to strictly positive weights simultaneously on
the countable edge set.

\begin{theorem}\label{thm:main}
Let $G$ be any nonatomic Borel probability measure on $[0,\infty)$, and let
$(t_e)_{e\in\Ecal}$ be independent with common law $G$. Under the stipulated
source convention below, almost surely there is no bigeodesic for $T_t$.
The assertion holds simultaneously over all bi-infinite simple paths,
without a prescribed direction or any moment assumption.
\end{theorem}

Throughout, \textbf{S} denotes OpenAI's \emph{No bigeodesics in planar
first-passage percolation}, September 24, 2026, at immutable revision
\texttt{fd4aeeb2ee4fc729c18d98444fed42fd0529eeeb} of the OpenAI math
repository~\cite{S}. Its theorem and its proof inputs under their stated
hypotheses are stipulated mathematical inputs to this deduction. We do not
audit that upstream theorem and do not apply its iid conclusion to a
dependent field. We use its actual geometric, measurable-selection, and
counting proof operations only after supplying the interfaces below.
\textbf{AH} means Ahlberg--Hoffman, arXiv:1609.02447v2~\cite{AH};
\textbf{DH} means Damron--Hanson, arXiv:1209.3036v1~\cite{DH}.
Their primary results are invoked with the hypothesis distinctions in
Section~\ref{sec:uses}.

S, Theorem 1.1, already covers \emph{every} nonatomic $G$ for which
\begin{equation}\label{eq:paid}
 \EE\big[\min(X_1,X_2,X_3,X_4)^2\big]<\infty,
 \qquad X_i\text{ independent with law }G.
\end{equation}
This includes bounded laws, unbounded finite-mean laws, and some
infinite-mean laws. The additional class in Theorem~\ref{thm:main} is exactly
the class where the expectation in \eqref{eq:paid} is infinite. We may
therefore assume below that $G$ is unbounded. No positive moment is then
assumed. The full stated scope agrees with the finite-valued formulation of
Damron--Hanson's Conjecture 1 in arXiv:1512.00804v2~\cite{DHb}.
This comparison identifies the scope; the deduction here remains conditional
on the stipulated input S.

\paragraph{Why an integrable surrogate is needed.}
The minimum moment in \eqref{eq:paid} is necessary for the usual almost-sure
uniform shape theorem for the original metric. We use the account of
Kesten's necessity and Cox--Durrett/Kesten contour regularization in
Cerf--Th\'eret~\cite[Section 1.2]{CT}; the original Cox--Durrett and Kesten
bodies have not been read for this manuscript. Thus no uniform original
shape or original Busemann linearity is asserted for the additional class.
The cap below has metric regularity, but fails upward finite energy
(Proposition~\ref{prop:energy}). Neither metric is an immediate black-box
input to all the structural results needed by S.

\paragraph{What is new in this deduction.}
The replacement steps are the dominated lifted limiting-graph coupling
(D, Section~\ref{sec:reference}), label and selector transfer under the
underlying-coordinate coupling (S, Section~\ref{sec:labels}), the changed-vertex
half-plane transfer (P, Section~\ref{sec:splice}), and the short cubic wall
(W, Section~\ref{sec:wall}). The circuit cap, exact bigeodesic equivalence,
shape estimates, same-status lowering, two-status density formulas, and S's
budget/fork/tilt argument receive no separate novelty claim.

Open-circuit surgery and its moment-free use have classical antecedents:
Auffinger--Damron--Hanson~\cite[Appendix C, Lemma C.5 and (26)]{ADH} use
open half-circuits to bypass costly geodesic portions; Cox--Durrett/Kesten
contours regularize passage times as described in~\cite{CT}. The direct
wall antecedent is AH, Proposition 9.10, Case 1. The stationary shape theorem
through Boivin is credited through AH Section 3.1 and DH's A2 discussion;
Boivin's original proof is not claimed read. Section~\ref{sec:metric}
retains an elementary shape derivation for completeness, not as a new result.

\paragraph{Logical order.}
First build the cap and its metric estimates. Construct coalescing reference
families by the modified DH argument. Reuse AH's order/label proofs, replacing
the downward step and then the upward selector step. This yields the all-ray
structure. Establish the independent-splice interface, then reuse S's
budget and divergent-toll contradiction with status-preserving kernels.
Every geometric bigeodesic is unchanged by the cap, so the conclusion returns
to the original metric. Earlier clipping, linear-cap, retention, and penalized-path
screens are historical material only and are not premises of this proof.
